\chapter{Cross-Cloud Service Equivalence}
\index{cross-cloud equivalence}\index{service mapping}\index{multi-cloud}%
This appendix is a living reference that maps conceptually equivalent services and
ideas across the seven platforms studied in this book series: Amazon Web Services
(AWS), Microsoft Azure, Microsoft Fabric, Google Cloud Platform (GCP), MongoDB,
Snowflake, and Databricks. The names differ, but the underlying data-engineering
concepts---durable object storage, tiered lifecycle economics, immutability, managed
relational engines, elastic compute, governed lakehouse tables, and disciplined data
organization---recur everywhere.

\begin{notebox}
This appendix grows with each chapter. Every chapter contributes the equivalence
table introduced in its cross-cloud callout, so by the end of the book you hold a
single consolidated Rosetta Stone for moving fluently between clouds. Databricks
appears throughout as the lakehouse-platform column: it runs on AWS, Azure, and GCP
alike, so its entries describe the platform layer rather than the underlying cloud.
Where a clean one-to-one mapping does not exist, the prose notes the closest analog
and the caveat.
\end{notebox}

\section{Free Tiers and Environment Setup}
\index{free tier!equivalence}\index{environment setup!equivalence}%
Chapter~0's account-and-tooling setup has a counterpart on every platform, and the
guardrail pattern---separate the admin identity, set a budget on day one, name and
tag everything for deletion---is identical everywhere.

\begin{center}
\footnotesize
\begin{tabular}{@{}p{2.9cm}p{3.0cm}p{3.0cm}p{3.0cm}@{}}
\toprule
\textbf{Setup step} & \textbf{AWS} & \textbf{Azure} & \textbf{GCP} \\
\midrule
Free offer & Free Tier (12 mo + always free) & Free account (\$200/30 days + always free) & Free trial (\$300/90 days + Always Free) \\
Identity model & IAM / IAM Identity Center & Microsoft Entra ID + RBAC & Cloud IAM \\
CLI & AWS CLI v2 & Azure CLI (\texttt{az login}) & gcloud CLI (\texttt{gcloud init}) \\
Browser shell & AWS CloudShell & Azure Cloud Shell & Cloud Shell \\
Cost guardrail & AWS Budgets + alerts & Cost Management budgets & Budgets \& alerts \\
\bottomrule
\end{tabular}
\end{center}

\noindent \textbf{Microsoft Fabric} starts from a Microsoft 365/Power BI work account with a trial
capacity; \textbf{Snowflake} offers a 30-day/\$400 trial guarded by Resource Monitors;
\textbf{MongoDB Atlas} has a forever-free M0 cluster behind database users and an IP allowlist.
\textbf{Databricks} offers a time-limited free trial on each cloud (plus the serverless free
tier and the Community Edition for learning); the first-day disciplines are the same---create a
workspace in its own resource group, attach cluster policies before users arrive, and route
billing through tags or \texttt{system.billing.usage} from the start.

\section{Object and Data-Lake Storage}
\index{object storage!equivalence}\index{storage classes!equivalence}%
The foundational layer covered in Chapter~1. Every platform offers durable, HTTP-addressable
object storage with tiered pricing, lifecycle transitions, versioning, and write-once-read-many
(WORM) controls.

\begin{center}
\footnotesize
\begin{tabular}{@{}p{2.8cm}p{3.2cm}p{3.2cm}p{3.2cm}@{}}
\toprule
\textbf{Capability} & \textbf{AWS} & \textbf{Azure} & \textbf{GCP} \\
\midrule
Object storage & Amazon S3 & ADLS Gen2 / Blob & Cloud Storage \\
Hot tier & S3 Standard & Hot & Standard \\
Infrequent access & S3 Standard-IA & Cool & Nearline \\
Archive (instant) & Glacier Instant Retrieval & Cold & Coldline \\
Deep archive & Glacier Deep Archive & Archive & Archive \\
Lifecycle rules & Lifecycle policies & Lifecycle management & Object Lifecycle Mgmt \\
Versioning & S3 Versioning & Blob versioning & Object Versioning \\
Immutability (WORM) & Object Lock & Immutable blob storage & Bucket Lock \\
Cross-region copy & CRR / SRR & Object replication & Dual/multi-region + Transfer \\
Encryption at rest & SSE-S3 / SSE-KMS & SSE + CMK (Key Vault) & Google-managed / CMEK \\
\bottomrule
\end{tabular}
\end{center}

\noindent The three remaining platforms approach storage differently.
\textbf{Microsoft Fabric} exposes a single logical lake, \emph{OneLake}, and uses
\emph{shortcuts} to virtualize S3, ADLS Gen2, and Google Cloud Storage in place without copying.
\textbf{Snowflake} abstracts storage entirely: table data lives in cloud object storage as
immutable, columnar \emph{micro-partitions} that the service manages on your behalf.
\textbf{MongoDB} persists documents through the \emph{WiredTiger} storage engine, balancing an
internal cache against the operating-system page cache rather than exposing storage tiers.
\textbf{Databricks} does not own a storage tier at all: it reads and writes the customer's S3,
ADLS Gen2, or GCS buckets through Unity Catalog \emph{external locations} and \emph{volumes},
with Delta Lake supplying the transactional table layer on top---proof that the lake and the
compute plane are separable purchases.

\section{Managed Relational Databases}
\index{relational databases!equivalence}%
The focus of Chapter~2 for the three hyperscalers. Each provides a fully managed relational
engine, a cloud-native high-performance variant, high availability, read replicas, and a
migration service.

\begin{center}
\footnotesize
\begin{tabular}{@{}p{2.8cm}p{3.2cm}p{3.2cm}p{3.2cm}@{}}
\toprule
\textbf{Capability} & \textbf{AWS} & \textbf{Azure} & \textbf{GCP} \\
\midrule
Managed relational & Amazon RDS & Azure SQL Database & Cloud SQL \\
Cloud-native engine & Amazon Aurora & SQL Hyperscale & AlloyDB \\
Managed PostgreSQL & RDS / Aurora PostgreSQL & Azure DB for PostgreSQL & Cloud SQL / AlloyDB \\
High availability & Multi-AZ & Zone-redundant HA & Regional HA \\
Read scaling & Read Replicas & Read replicas / geo-replicas & Read Replicas \\
Global / geo & Aurora Global Database & Active geo-replication & Cross-region replicas \\
Serverless compute & Aurora Serverless v2 & Azure SQL Serverless & AlloyDB (elastic) \\
Migration service & DMS + SCT & Azure Database Migration Svc & Database Migration Svc \\
\bottomrule
\end{tabular}
\end{center}

\noindent \textbf{Snowflake}, \textbf{Fabric}, and \textbf{MongoDB} are not managed relational
engines, but they solve overlapping problems: Snowflake and Fabric provide SQL analytics over
warehouse/lakehouse storage, while MongoDB provides a distributed document store whose schema
design (Chapter~2) replaces rigid relational normalization with intentional, access-driven modeling.
\textbf{Databricks} is likewise not an OLTP engine; its \emph{Lakehouse Federation} feature queries
external relational systems (SQL Server, PostgreSQL, MySQL) in place so operational data can be
joined or ingested without a separate ETL hop.

\section{Elastic Compute and Analytical Warehouses}
\index{elastic compute!equivalence}\index{virtual warehouse!equivalence}%
Snowflake's Chapter~2 topic---separating compute from storage and scaling it elastically---has a
counterpart in every analytics platform.

\begin{center}
\footnotesize
\begin{tabular}{@{}p{2.8cm}p{3.2cm}p{3.2cm}p{3.2cm}@{}}
\toprule
\textbf{Concept} & \textbf{Snowflake} & \textbf{AWS} & \textbf{Azure / GCP} \\
\midrule
Compute unit & Virtual Warehouse & Redshift RA3 / Serverless & Synapse pool / BigQuery slots \\
Elastic concurrency & Multi-cluster warehouse & Concurrency scaling & Autoscale / on-demand \\
Pause when idle & Auto-suspend / resume & Pause / resume & Auto-pause / serverless \\
Billing unit & Credits & RPU / node-hours & DWU / slot-hours \\
Result reuse & Result cache & Result cache & Result / BI cache \\
\bottomrule
\end{tabular}
\end{center}

\noindent \textbf{Microsoft Fabric} meters all workloads against a shared pool of Capacity Units
(CU) on an F-SKU, an approach closer to a single elastic budget than to per-engine warehouses.
\textbf{MongoDB Atlas} scales through cluster tiers with optional auto-scaling of compute and storage.
\textbf{Databricks} bills Databricks Units (DBU) on top of VM cost: all-purpose clusters for
interactive work, job clusters that terminate on completion for pipelines, and serverless SQL
warehouses that scale to zero---the same pause-when-idle discipline with more explicit knobs.

\section{Data Organization and the Medallion Pattern}
\index{medallion architecture!equivalence}%
Fabric's Chapter~2 topic---the Bronze/Silver/Gold medallion---is a portable design pattern rather
than a product. Every platform implements the same progressive-refinement idea with its own primitives.

\begin{center}
\footnotesize
\begin{tabular}{@{}p{2.6cm}p{3.4cm}p{3.4cm}p{3.0cm}@{}}
\toprule
\textbf{Layer} & \textbf{Fabric} & \textbf{Data lake (AWS/Azure/GCP)} & \textbf{Snowflake} \\
\midrule
Bronze (raw) & Bronze Lakehouse & \texttt{raw/} prefix (S3/ADLS/GCS) & RAW schema / stage \\
Silver (validated) & Silver Lakehouse & \texttt{validated/} prefix & STAGING schema \\
Gold (curated) & Gold Lakehouse & \texttt{curated/} prefix & ANALYTICS schema \\
Storage format & Delta-Parquet & Parquet / Delta / Iceberg & Micro-partitions \\
\bottomrule
\end{tabular}
\end{center}

\noindent \textbf{MongoDB} realizes the same separation with distinct collections or databases
(for example \texttt{raw\_events}, \texttt{validated}, \texttt{curated}) governed by schema-design
patterns rather than a lake layout. \textbf{Databricks} is where the medallion vocabulary was
popularized: \texttt{bronze}/\texttt{silver}/\texttt{gold} schemas of Delta tables under a Unity
Catalog catalog, with expectations and quarantine as the quality gates between layers and Unity
Catalog grants as the access boundary per layer.

\section{Document Data Modeling}
\index{document modeling!equivalence}%
MongoDB's Chapter~2 schema-design patterns generalize to every document and wide-column store.

\begin{center}
\footnotesize
\begin{tabular}{@{}p{3.0cm}p{3.2cm}p{3.2cm}p{3.0cm}@{}}
\toprule
\textbf{Concept} & \textbf{MongoDB} & \textbf{Amazon DynamoDB} & \textbf{Azure Cosmos DB} \\
\midrule
Primary access unit & Document / collection & Item / table & Item / container \\
Partitioning & Shard key & Partition key & Partition key \\
One-table modeling & Embedding + patterns & Single-table design & Embedding + partitioning \\
Time-series grouping & Bucket Pattern & Composite sort keys & Bucketing by key \\
Optional attributes & Attribute Pattern & Sparse attributes & Sparse properties \\
\bottomrule
\end{tabular}
\end{center}

\noindent Google Cloud's closest analogs are \textbf{Firestore} and \textbf{Bigtable}. In
\textbf{Snowflake} and \textbf{Fabric}, semi-structured documents are stored in \texttt{VARIANT}
(Snowflake) or JSON columns and queried with path expressions rather than modeled as collections.
\textbf{Databricks} follows the same lakehouse pattern: JSON payloads land as strings or
\texttt{VARIANT} columns in Delta tables and are shredded with \texttt{from\_json} and schema
hints during the bronze-to-silver pass.

\section{Globally Distributed and NoSQL Databases}
\index{NoSQL!equivalence}\index{distributed databases!equivalence}%
Chapter~3 branches by platform into high-scale NoSQL (AWS DynamoDB, Azure Cosmos DB), globally
consistent NewSQL (GCP Cloud Spanner), and advanced document modeling (MongoDB). They all answer
the same question: how to scale a database horizontally across partitions and regions.

\begin{center}
\footnotesize
\begin{tabular}{@{}p{2.8cm}p{2.9cm}p{2.9cm}p{2.9cm}p{2.4cm}@{}}
\toprule
\textbf{Concept} & \textbf{DynamoDB} & \textbf{Cosmos DB} & \textbf{Cloud Spanner} & \textbf{MongoDB} \\
\midrule
Model & Key-value / doc & Multi-model & Distributed SQL & Document \\
Partitioning & Partition key & Partition key & Split by key range & Shard key \\
Throughput unit & RCU / WCU & Request Units (RU) & Nodes / processing units & Cluster tier \\
Secondary indexes & GSI / LSI & Automatic indexing & Secondary indexes & Indexes \\
Change stream & DynamoDB Streams & Change Feed & Change streams & Change streams \\
Expiry & TTL & TTL & (manual) & TTL index \\
Multi-region & Global Tables & Multi-region writes & Multi-region (synchronous) & Zones / global clusters \\
Consistency & Eventual / strong & 5 tunable levels & External (TrueTime) & Tunable read/write concern \\
\bottomrule
\end{tabular}
\end{center}

\noindent \textbf{Snowflake} and \textbf{Fabric} are analytical platforms rather than OLTP key-value
stores; they ingest from these operational databases and store semi-structured payloads in
\texttt{VARIANT}/JSON columns. Google's dedicated NoSQL analogs are \textbf{Firestore} (document)
and \textbf{Bigtable} (wide-column). \textbf{Databricks} deliberately stays out of operational
NoSQL: it consumes change feeds from Cosmos DB, DynamoDB, and MongoDB into Delta tables, leaving
point reads and sub-millisecond serving to the operational stores.

\section{Managed Apache Spark}
\index{Apache Spark!equivalence}%
Microsoft Fabric's Chapter~3 topic---managed Spark pools and notebooks---has a first-class
equivalent on every cloud.

\begin{center}
\footnotesize
\begin{tabular}{@{}p{2.8cm}p{3.2cm}p{3.2cm}p{3.2cm}@{}}
\toprule
\textbf{Capability} & \textbf{Fabric} & \textbf{AWS} & \textbf{Azure / GCP} \\
\midrule
Managed Spark & Starter / Custom pools & Glue / EMR / EMR Serverless & Synapse Spark / Databricks; Dataproc \\
Notebook UX & Fabric Notebooks & Glue Studio / EMR Studio & Synapse / Databricks; Vertex AI \\
Serverless option & Autoscaling pools & Glue / EMR Serverless & Synapse serverless; Dataproc Serverless \\
Lake format & Delta on OneLake & Delta / Iceberg / Hudi & Delta; Delta / Iceberg \\
\bottomrule
\end{tabular}
\end{center}

\noindent \textbf{Snowflake} offers \textbf{Snowpark} for DataFrame-style processing, and
\textbf{MongoDB} integrates through the Spark Connector. \textbf{Databricks} is the reference
implementation this book uses in Chapter~9: the Databricks Runtime with the vectorized Photon
engine, a built-in cluster manager (no YARN to run), autoscaling clusters plus serverless
compute, and Delta as the default table format. The Spark API, Catalyst optimizer, and shuffle
mechanics are identical on EMR, Dataproc, Fabric, and Databricks---only the operational wrapper
and the billing unit change.

\section{Query Acceleration and Materialized Views}
\index{materialized views!equivalence}\index{query acceleration!equivalence}%
Snowflake's Chapter~3 acceleration tools---materialized views, the Search Optimization Service, and
transient/temporary tables---map to comparable features in every warehouse.

\begin{center}
\footnotesize
\begin{tabular}{@{}p{2.8cm}p{3.0cm}p{3.0cm}p{3.4cm}@{}}
\toprule
\textbf{Feature} & \textbf{Snowflake} & \textbf{AWS Redshift} & \textbf{Azure Synapse / GCP BigQuery} \\
\midrule
Materialized views & Materialized Views & Materialized views & Materialized views \\
Point-lookup accel. & Search Optimization Svc & Sort keys / zone maps & (indexing) / search indexes \\
Result reuse & Result cache & Result cache & Result-set cache / BI Engine \\
Ephemeral tables & Transient / temporary & Temp tables & Temp tables \\
\bottomrule
\end{tabular}
\end{center}

\noindent \textbf{Fabric} accelerates reads through Direct Lake and its Warehouse engine, while
\textbf{MongoDB} relies on secondary and compound indexes plus the in-memory working set.
\textbf{Databricks SQL} pairs materialized views and streaming tables with a result cache on
serverless warehouses, and adds file-level acceleration---\texttt{OPTIMIZE}, Z-order and liquid
clustering, bloom filters, and deletion vectors---that the next sections cover in detail.

\section{In-Memory Caching and Hybrid Ingestion}
\index{caching!equivalence}\index{in-memory cache!equivalence}%
The AWS and Azure Chapter~4 topic---low-latency caching in front of a durable store, fed by hybrid
(edge-to-cloud) ingestion---has a managed Redis-compatible equivalent on every cloud.

\begin{center}
\footnotesize
\begin{tabular}{@{}p{2.8cm}p{3.0cm}p{3.0cm}p{3.4cm}@{}}
\toprule
\textbf{Capability} & \textbf{AWS} & \textbf{Azure} & \textbf{GCP / others} \\
\midrule
Managed cache & ElastiCache (Redis/Memcached) & Cache for Redis & Memorystore (Redis/Memcached) \\
Real-time API layer & AppSync / API Gateway & API Management / SignalR & API Gateway / Apigee \\
Hybrid ingestion & DataSync / Snow family / IoT Core & Data Box / IoT Hub / Arc & Transfer Appliance / IoT Core \\
Streaming intake & Kinesis / MSK & Event Hubs & Pub/Sub \\
\bottomrule
\end{tabular}
\end{center}

\noindent \textbf{Snowflake} and \textbf{Fabric} rely on result and Direct~Lake caches rather than a
standalone Redis tier, while \textbf{MongoDB} serves hot data from its in-memory working set and
Atlas can pair with an external cache. \textbf{Databricks} caches at a different layer: the
disk cache (formerly Delta cache) keeps hot Parquet blocks on cluster-local SSDs, and
predictive I/O accelerates repeated scans without any application-managed cache keys.

\section{Wide-Column and Key-Value NoSQL}
\index{wide-column!equivalence}\index{Bigtable!equivalence}%
GCP's Chapter~4 topic---Cloud Bigtable---is a wide-column, high-throughput, low-latency NoSQL store.
Its peers span the managed and open-source ecosystems.

\begin{center}
\footnotesize
\begin{tabular}{@{}p{2.8cm}p{3.0cm}p{3.0cm}p{3.4cm}@{}}
\toprule
\textbf{Characteristic} & \textbf{GCP Bigtable} & \textbf{AWS} & \textbf{Azure / open source} \\
\midrule
Data model & Wide-column & DynamoDB (key-value/doc) & Cosmos DB (multi-model) \\
HBase-compatible & HBase API & Keyspaces (Cassandra) & Managed Instance for Cassandra \\
Scaling unit & Nodes / clusters & On-demand / provisioned & RU/s (Cosmos) \\
Best fit & Time-series, IoT, ad-tech & Serverless key-value & Global multi-model \\
\bottomrule
\end{tabular}
\end{center}

\noindent \textbf{Snowflake} and \textbf{Fabric} address these workloads analytically rather than as
operational key-value stores; \textbf{MongoDB} covers the document end of the same NoSQL spectrum.
\textbf{Databricks} likewise serves these workloads analytically: Bigtable and Cassandra estates
are typically migrated by bulk export to Delta tables, not by replacing the serving path.

\section{Data Protection: Cloning, Time Travel, and Migration}
\index{cloning!equivalence}\index{time travel!equivalence}\index{point-in-time recovery!equivalence}%
Snowflake's Chapter~4 topic---zero-copy cloning, Time Travel, and Fail-safe---plus MongoDB's
relational-to-document migration map to data-protection and migration tooling across the clouds.

\begin{center}
\footnotesize
\begin{tabular}{@{}p{2.8cm}p{3.0cm}p{3.0cm}p{3.4cm}@{}}
\toprule
\textbf{Capability} & \textbf{Snowflake} & \textbf{AWS / Azure} & \textbf{GCP / Fabric} \\
\midrule
Instant clone & Zero-copy clone & Aurora clone / SQL DB copy & BigQuery table clone; OneLake shortcut \\
Point-in-time recovery & Time Travel & PITR (RDS/SQL DB backups) & PITR; Delta time travel \\
Retention safety net & Fail-safe (7 days) & Automated backups & Backups; Delta VACUUM window \\
Schema migration & --- & DMS / SCT & Database Migration Service \\
\bottomrule
\end{tabular}
\end{center}

\noindent \textbf{MongoDB} migrations use \textbf{Relational Migrator} plus the embed-vs-reference
patterns from Chapter~4, and \textbf{Delta Lake} exposes historical versions through
\texttt{VERSION AS OF} time-travel queries. \textbf{Databricks} operationalizes the same
protection set on Delta: \texttt{RESTORE} for point-in-time recovery, \texttt{VACUUM} as the
retention safety-valve cutoff, and \emph{shallow clones} as the zero-copy clone analog for
dev/test copies of production tables.

\section{MPP Data Warehouses}
\index{MPP warehouse!equivalence}\index{Redshift!equivalence}%
Chapter~5's dedicated SQL pool is one of the most portable ideas in analytics engineering: a
control node coordinates compute nodes over physically distributed tables. The same mental model
applies to Redshift's ra3 clusters and BigQuery's slot-based execution.

\begin{center}
\footnotesize
\begin{tabular}{@{}p{3.0cm}p{2.9cm}p{3.1cm}p{3.2cm}@{}}
\toprule
\textbf{Capability} & \textbf{AWS} & \textbf{Azure} & \textbf{GCP} \\
\midrule
MPP warehouse & Amazon Redshift & Synapse dedicated SQL pool & BigQuery (slot-based) \\
Sizing unit & ra3 nodes / RPUs & cDWU (pauseable) & Slots / autoscaling editions \\
Distribution & DISTKEY / DISTSTYLE & Hash / round-robin / replicated & Automatic (no user control) \\
Bulk load from lake & COPY from S3 & COPY INTO from ADLS Gen2 & LOAD DATA from Cloud Storage \\
Serverless SQL over lake & Athena / Redshift Spectrum & Synapse serverless SQL pool & BigQuery external tables \\
\addlinespace
\multicolumn{4}{@{}l@{}}{\textbf{Fabric}: T-SQL MPP warehouse on OneLake; \textbf{Snowflake}: virtual warehouses; \textbf{MongoDB}: Atlas SQL (not MPP);} \\
\multicolumn{4}{@{}l@{}}{\textbf{Databricks}: SQL warehouses (serverless / pro / classic) with the Photon engine over Delta tables.} \\
\bottomrule
\end{tabular}
\end{center}

\noindent The Kimball discipline---grain declarations, conformed dimensions, slowly changing
dimensions---is engine-neutral. A star schema designed well in Chapter~5 ports to Redshift,
BigQuery, Snowflake, a Fabric warehouse, or a Databricks SQL warehouse with only dialect changes.

\section{Warehouse Physical Design and Workload Isolation}
\index{workload isolation!equivalence}\index{columnstore!equivalence}%
Chapter~6's optimization levers---columnar storage, distribution, workload groups, precomputed
results---exist on every warehouse; only the knobs differ. Synapse exposes physical design
explicitly, while BigQuery and Snowflake automate more of it: the trade-off is control versus toil.

\begin{center}
\footnotesize
\begin{tabular}{@{}p{3.0cm}p{2.9cm}p{3.1cm}p{3.2cm}@{}}
\toprule
\textbf{Capability} & \textbf{AWS (Redshift)} & \textbf{Azure (Synapse)} & \textbf{GCP (BigQuery)} \\
\midrule
Columnar storage & Columnar by default & Clustered columnstore index & Columnar by default \\
Physical layout & Sort keys + dist keys & Distribution + partitioning & Clustering + partitioning \\
Workload isolation & WLM queues & Workload groups + classifiers & Slot reservations \\
Precomputed results & Materialized views & Materialized views + result-set cache & Materialized views + BI Engine \\
Query cache & Result cache & Result-set cache & Cached query results \\
\addlinespace
\multicolumn{4}{@{}l@{}}{\textbf{Fabric}: warehouse with automatic tuning; \textbf{Snowflake}: micro-partition pruning + resource monitors;} \\
\multicolumn{4}{@{}l@{}}{\textbf{MongoDB}: indexes + covered queries; \textbf{Databricks}: \texttt{OPTIMIZE}, Z-order, liquid clustering, predictive optimization.} \\
\bottomrule
\end{tabular}
\end{center}

\noindent Databricks' \emph{liquid clustering} deserves special attention: like Iceberg's sort-order
evolution and BigQuery's clustering, it makes physical layout rewritable without rewriting the
table, which removes the classic ``wrong distribution key'' lock-in that Chapter~6 warns about.

\section{Serverless SQL over the Lake}
\index{serverless SQL!equivalence}\index{Athena!equivalence}%
Chapter~7's pattern---query the lake in place, pay per byte scanned, materialize back with
CETAS---exists on every major cloud, and the billing model makes format and partitioning
discipline a cross-cloud skill rather than a Synapse trick.

\begin{center}
\footnotesize
\begin{tabular}{@{}p{3.0cm}p{2.9cm}p{3.1cm}p{3.2cm}@{}}
\toprule
\textbf{Capability} & \textbf{AWS} & \textbf{Azure} & \textbf{GCP} \\
\midrule
Serverless lake SQL & Amazon Athena & Synapse serverless SQL pool & BigQuery external tables \\
Query-in-place API & S3 Select / Athena & OPENROWSET & External data sources \\
Materialize to lake & CTAS into S3 & CETAS & Export / CREATE TABLE AS \\
Open table format & Delta / Iceberg / Hudi on S3 & Delta Lake on ADLS Gen2 & BigLake tables / Iceberg \\
Billing lever & Per TB scanned & Per TB scanned & Per TB scanned (on-demand) \\
\addlinespace
\multicolumn{4}{@{}l@{}}{\textbf{Fabric}: SQL analytics endpoint over OneLake; \textbf{Snowflake}: external tables / Iceberg;} \\
\multicolumn{4}{@{}l@{}}{\textbf{MongoDB}: Atlas Data Federation over S3; \textbf{Databricks}: serverless SQL warehouses over UC Delta tables.} \\
\bottomrule
\end{tabular}
\end{center}

\noindent The lakehouse thesis---warehouse-class reliability on lake storage---originated with
Delta Lake \cite{armbrust2020delta,armbrust2021lakehouse} and is now the default architecture
vocabulary across all seven platforms in this series.

\section{Data Catalogs and Governance Planes}
\index{catalog!equivalence}\index{Purview!equivalence}%
Chapter~8's governance plane---a metastore-backed catalog, automated classification, lineage, and
in-place sharing---has a native equivalent on every cloud. Purview's distinction is scope: it
catalogs Azure, AWS, GCP, and on-premises sources in one data map.

\begin{center}
\footnotesize
\begin{tabular}{@{}p{3.0cm}p{2.9cm}p{3.1cm}p{3.2cm}@{}}
\toprule
\textbf{Capability} & \textbf{AWS} & \textbf{Azure} & \textbf{GCP} \\
\midrule
Catalog / data map & Glue Data Catalog + Lake Formation & Microsoft Purview & Dataplex \\
PII discovery & Macie & Purview classifications & Sensitive Data Protection \\
Lineage & Glue / CloudTrail analysis & Purview lineage & Dataplex lineage \\
Cross-tenant sharing & AWS Data Exchange / S3 sharing & Azure Data Share & Analytics Hub / BigQuery sharing \\
Clean room & AWS Clean Rooms & Data Share + confidential compute & BigQuery data clean rooms \\
\addlinespace
\multicolumn{4}{@{}l@{}}{\textbf{Fabric}: Purview hub + OneLake catalog; \textbf{Snowflake}: Horizon Catalog + listings;} \\
\multicolumn{4}{@{}l@{}}{\textbf{MongoDB}: Atlas auditing; \textbf{Databricks}: Unity Catalog + Delta Sharing (open protocol).} \\
\bottomrule
\end{tabular}
\end{center}

\noindent Unity Catalog's distinctive move is one governance model---grants, row filters, column
masks, lineage, audit---spanning tables, files, notebooks, models, and functions; the
hyperscalers split the same job across a catalog plus a policy engine. Delta Sharing's
distinction is the open protocol: recipients consume shares with pandas, Spark, or Power BI
without any Databricks account.

\section{Managed Visual ETL and Orchestration}
\index{ETL service!equivalence}\index{AWS Glue!equivalence}%
Chapter~10's four-part ADF object model---connection (linked service), data pointer (dataset),
activity graph (pipeline), schedule (trigger)---has a direct analog in every managed ETL service.

\begin{center}
\footnotesize
\begin{tabular}{@{}p{3.0cm}p{2.9cm}p{3.1cm}p{3.2cm}@{}}
\toprule
\textbf{Capability} & \textbf{AWS} & \textbf{Azure} & \textbf{GCP} \\
\midrule
Managed ETL service & AWS Glue & Azure Data Factory & Cloud Data Fusion \\
Visual transform engine & Glue Studio / DataBrew & Mapping Data Flows & Data Fusion wrangler \\
Orchestration & Step Functions / MWAA & ADF pipelines and triggers & Cloud Composer (Airflow) \\
Serverless transform compute & Glue jobs (Spark) & Data Flows on managed Spark & Dataflow (Beam) \\
Connection abstraction & Glue connections & Linked services & Data Fusion connections \\
\addlinespace
\multicolumn{4}{@{}l@{}}{\textbf{Fabric}: pipelines + dataflows gen2; \textbf{Snowflake}: tasks/streams + partner tools;} \\
\multicolumn{4}{@{}l@{}}{\textbf{MongoDB}: Atlas Triggers; \textbf{Databricks}: Lakeflow Declarative Pipelines + Workflows.} \\
\bottomrule
\end{tabular}
\end{center}

\noindent Mapping Data Flows compile to Spark under the hood, so the execution concepts from
Chapter~9 (partitions, shuffles) apply underneath the visual canvas---the abstraction leaks
upward when performance matters. Databricks' declarative pipelines share the same philosophy:
define datasets and quality expectations, let the platform handle execution, retries, and
incremental refresh.

\section{Change Data Capture, Migration, and Hybrid Connectivity}
\index{CDC!equivalence}\index{data migration!equivalence}%
Chapter~11's integration problems---reaching private networks and streaming database changes---are
universal. Every cloud offers an agent/relay for private connectivity and a replication service
for change data capture.

\begin{center}
\footnotesize
\begin{tabular}{@{}p{3.0cm}p{2.9cm}p{3.1cm}p{3.2cm}@{}}
\toprule
\textbf{Capability} & \textbf{AWS} & \textbf{Azure} & \textbf{GCP} \\
\midrule
Private data movement agent & Glue connectors / DataSync agent & Self-Hosted Integration Runtime & VPN/Interconnect + remote agent \\
Database migration & AWS DMS (+SCT) & Azure DMS (+DMA) & Database Migration Service \\
CDC stream & DMS ongoing replication & Azure SQL CDC / DMS CDC & Datastream \\
Event fabric & EventBridge & Event Grid & Eventarc \\
Schema enforcement & Glue Schema Registry & Contract checks in ADF/Databricks & Dataflow / Dataplex validation \\
\addlinespace
\multicolumn{4}{@{}l@{}}{\textbf{Fabric}: on-premises data gateway; \textbf{Snowflake}: Snowpipe + connectors;} \\
\multicolumn{4}{@{}l@{}}{\textbf{MongoDB}: change streams (CDC-native); \textbf{Databricks}: Lakeflow Connect, Auto Loader, \texttt{AUTO CDC}.} \\
\bottomrule
\end{tabular}
\end{center}

\noindent Data contracts---schema, versioning, quarantine---are platform-neutral governance.
Databricks encodes them directly in the pipeline: Auto Loader's schema inference with hints and
evolution, the \texttt{\_rescued\_data} column for malformed rows, and DLT expectations that warn,
drop, or fail on contract violations.

\section{Medallion Orchestration and Scheduling}
\index{orchestration!equivalence}\index{Airflow!equivalence}%
Chapter~12's medallion pipelines need a scheduler with dependency graphs and failure semantics on
every platform. Bronze/silver/gold is the same idea whether the lake is S3, ADLS, or GCS.

\begin{center}
\footnotesize
\begin{tabular}{@{}p{3.0cm}p{2.9cm}p{3.1cm}p{3.2cm}@{}}
\toprule
\textbf{Capability} & \textbf{AWS} & \textbf{Azure} & \textbf{GCP} \\
\midrule
Medallion zones & S3 raw/curated + Glue & ADLS Gen2 + Delta Lake & GCS + BigLake \\
Managed Airflow & MWAA & Managed Airflow in ADF & Cloud Composer \\
Control-flow orchestration & Step Functions & ADF pipeline activities & Workflows \\
Event trigger & EventBridge $\rightarrow$ Lambda/Glue & Event Grid $\rightarrow$ ADF trigger & Eventarc $\rightarrow$ Workflows \\
Notification workflows & SNS / Step Functions & Logic Apps & Cloud Functions + Workflows \\
\addlinespace
\multicolumn{4}{@{}l@{}}{\textbf{Fabric}: OneLake medallion + pipelines; \textbf{Snowflake}: tasks/streams DAGs;} \\
\multicolumn{4}{@{}l@{}}{\textbf{MongoDB}: Atlas Triggers; \textbf{Databricks}: Workflows with repair-run partial recovery + job clusters.} \\
\bottomrule
\end{tabular}
\end{center}

\noindent Airflow's DAG model---tasks, dependencies, retries, backfills---is itself the cross-cloud
skill: MWAA, Cloud Composer, and Managed Airflow in ADF all execute the same Python DAG files
\cite{apache2025airflow}, and Databricks Workflows integrates with external orchestrators through
the Jobs API when an estate standardizes on Airflow.

\section{Event Streaming Backbones}
\index{event streaming!equivalence}\index{Kafka!equivalence}%
Chapter~13's partitioned, replayable event log is a cross-cloud primitive. Event Hubs'
model---partitions, consumer groups, offsets, retention---is deliberately Kafka-shaped, and every
cloud offers an equivalent.

\begin{center}
\footnotesize
\begin{tabular}{@{}p{3.0cm}p{2.9cm}p{3.1cm}p{3.2cm}@{}}
\toprule
\textbf{Capability} & \textbf{AWS} & \textbf{Azure} & \textbf{GCP} \\
\midrule
Event streaming & Kinesis Data Streams / MSK & Event Hubs (+ Kafka endpoint) & Pub/Sub / Managed Kafka \\
Message queue & SQS / Amazon MQ & Service Bus queues and topics & Pub/Sub (pull) \\
Device ingestion & IoT Core & IoT Hub & (partner; IoT Core retired) \\
Stream-to-lake export & Kinesis Firehose & Event Hubs Capture & Pub/Sub to GCS subscriptions \\
Consumer isolation & Enhanced fan-out & Consumer groups & Subscriptions \\
\addlinespace
\multicolumn{4}{@{}l@{}}{\textbf{Fabric}: eventstreams into OneLake; \textbf{Snowflake}: Snowpipe Streaming;} \\
\multicolumn{4}{@{}l@{}}{\textbf{MongoDB}: change streams as a source; \textbf{Databricks}: \texttt{kafka} format source against the Event Hubs Kafka endpoint.} \\
\bottomrule
\end{tabular}
\end{center}

\noindent Because Event Hubs speaks the Kafka protocol, producer and consumer skills transfer to
Amazon MSK, Confluent, and self-managed clusters with only endpoint and credential changes
\cite{apache2025kafka}---which is exactly how Chapter~15's Databricks consumer attaches.

\section{Managed SQL over Streams}
\index{stream SQL!equivalence}\index{windowing!equivalence}%
Chapter~14's managed SQL-over-streams exists on every cloud, and the windowing
vocabulary---tumbling, hopping, sliding, session---is identical everywhere; only the engine and
billing unit differ.

\begin{center}
\footnotesize
\begin{tabular}{@{}p{3.0cm}p{2.9cm}p{3.1cm}p{3.2cm}@{}}
\toprule
\textbf{Capability} & \textbf{AWS} & \textbf{Azure} & \textbf{GCP} \\
\midrule
Managed stream SQL & Managed Flink (Kinesis Data Analytics) & Azure Stream Analytics & Dataflow (Beam SQL) \\
Billing unit & Kinesis Processing Units & Streaming Units (SUs) & Dataflow worker resources \\
Event source & Kinesis / MSK & Event Hubs / IoT Hub & Pub/Sub \\
Late-data control & Watermarks in Flink & Late/out-of-order policies & Watermarks + allowed lateness \\
Stream to warehouse & Firehose $\rightarrow$ Redshift & ASA output $\rightarrow$ Synapse/SQL & Dataflow $\rightarrow$ BigQuery \\
\addlinespace
\multicolumn{4}{@{}l@{}}{\textbf{Fabric}: eventstreams + Activator; \textbf{Snowflake}: Snowpipe Streaming + dynamic tables;} \\
\multicolumn{4}{@{}l@{}}{\textbf{MongoDB}: Atlas Stream Processing; \textbf{Databricks}: Structured Streaming / DLT streaming tables.} \\
\bottomrule
\end{tabular}
\end{center}

\noindent When a use case outgrows ASA---complex state, stream-stream joins at scale, custom
code---Spark Structured Streaming (Chapter~15) is the escalation path, with the same windowing
concepts at a lower level of abstraction. Databricks exposes the full state-design surface:
\texttt{withWatermark}, \texttt{dropDuplicates}, and \texttt{flatMapGroupsWithState} with
RocksDB-backed state.

\section{Spark Structured Streaming}
\index{Structured Streaming!equivalence}%
Chapter~15's model---a stream as an unbounded table, micro-batches, watermarks,
checkpoints---maps directly onto the streaming engines of the other platforms.

\begin{center}
\footnotesize
\begin{tabular}{@{}p{3.0cm}p{2.9cm}p{3.1cm}p{3.2cm}@{}}
\toprule
\textbf{Capability} & \textbf{AWS} & \textbf{Azure} & \textbf{GCP} \\
\midrule
Spark streaming host & EMR / Glue streaming & Azure Databricks & Dataproc / Serverless Spark \\
Stream source & Kinesis / MSK connector & Event Hubs (Kafka endpoint) & Pub/Sub / Kafka connector \\
State store & RocksDB on EMR & Default or RocksDB provider & RocksDB on Dataproc \\
Exactly-once sink & Delta/Iceberg merge & Delta Lake merge & BigQuery Storage Write API \\
Checkpoint storage & S3 & ADLS Gen2 & GCS \\
\addlinespace
\multicolumn{4}{@{}l@{}}{\textbf{Fabric}: structured streaming on OneLake; \textbf{Snowflake}: Snowpipe Streaming + tasks;} \\
\multicolumn{4}{@{}l@{}}{\textbf{MongoDB}: Atlas Stream Processing; \textbf{Databricks}: native---the reference implementation used in this chapter.} \\
\bottomrule
\end{tabular}
\end{center}

\noindent The Structured Streaming programming guide \cite{apache2025structuredStreaming} is
platform-neutral reading; everything Chapter~15 teaches about watermarks and state applies
identically on EMR and Dataproc. The batch-unification trigger (\texttt{availableNow}) is the
same idea as Dataflow's unified batch/stream thesis and Flink's batch mode.

\section{Business Intelligence and Semantic Models}
\index{BI!equivalence}\index{semantic model!equivalence}%
Chapter~16's BI layer is the least portable part of the stack---each cloud has a native analytics
front end---but the underlying concepts (extracted versus live models, row-level security, measure
semantics) transfer directly.

\begin{center}
\footnotesize
\begin{tabular}{@{}p{3.0cm}p{2.9cm}p{3.1cm}p{3.2cm}@{}}
\toprule
\textbf{Capability} & \textbf{AWS} & \textbf{Azure} & \textbf{GCP} \\
\midrule
BI service & Amazon QuickSight & Power BI & Looker / Looker Studio \\
Live query mode & Direct query to Redshift/Athena & DirectQuery & Live connection to BigQuery \\
Extracted cache & SPICE & Import mode (in-memory) & Extracted data / BI Engine \\
Row-level security & QuickSight RLS & Power BI RLS roles & Looker access filters \\
Metrics layer & QuickSight datasets & Semantic model + DAX & LookML \\
\addlinespace
\multicolumn{4}{@{}l@{}}{\textbf{Fabric}: Power BI with DirectLake over OneLake; \textbf{Snowflake}: Snowsight + partner BI;} \\
\multicolumn{4}{@{}l@{}}{\textbf{MongoDB}: Atlas Charts + BI Connector; \textbf{Databricks}: AI/BI dashboards + Genie natural-language spaces.} \\
\bottomrule
\end{tabular}
\end{center}

\noindent Power BI's \textbf{DirectLake} mode in Microsoft Fabric is the emerging fourth storage
mode: the engine reads Delta tables in OneLake directly, combining Import-mode speed with
DirectQuery freshness. Databricks' AI/BI makes the mirror-image bet: dashboards read governed
Delta tables in place and inherit Unity Catalog lineage, so a tile traces back to the pipeline
that produced it.

\section{Feature Engineering and In-Database Machine Learning}
\index{feature store!equivalence}\index{in-database ML!equivalence}%
Chapter~17's in-database scoring---pushing the model to the data rather than shipping data to a
model server---exists in every analytical engine, typically via ONNX or a native PREDICT-style
function.

\begin{center}
\footnotesize
\begin{tabular}{@{}p{3.0cm}p{2.9cm}p{3.1cm}p{3.2cm}@{}}
\toprule
\textbf{Capability} & \textbf{AWS} & \textbf{Azure} & \textbf{GCP} \\
\midrule
In-database scoring & Redshift ML (SageMaker) & Synapse T-SQL PREDICT (ONNX) & BigQuery ML \\
Spark training & EMR + SageMaker & Synapse Spark / SynapseML & Dataproc + Vertex AI \\
Feature store & SageMaker Feature Store & AML managed feature store & Vertex AI Feature Store \\
Model format for SQL & (auto via SageMaker) & ONNX & Native BQML models \\
Batch scoring & Redshift batch transform & PREDICT over pool tables & ML.PREDICT \\
\addlinespace
\multicolumn{4}{@{}l@{}}{\textbf{Fabric}: SynapseML + MLflow in notebooks; \textbf{Snowflake}: Snowpark ML + SQL inference;} \\
\multicolumn{4}{@{}l@{}}{\textbf{MongoDB}: Atlas Vector Search scoring; \textbf{Databricks}: Feature Engineering in UC + Mosaic AI serving.} \\
\bottomrule
\end{tabular}
\end{center}

\noindent ONNX is the portable artifact: train in Spark, export once, score in SQL, in an AML
endpoint (Chapter~19), or at the edge. Databricks' feature store makes the anti-skew contract
explicit: one governed Delta definition feeds both training sets (via point-in-time joins) and
online serving lookups.

\section{Managed Machine Learning Platforms}
\index{ML platform!equivalence}\index{SageMaker!equivalence}%
Chapter~18's managed ML platforms share a universal anatomy: a workspace, versioned data
references, reusable environments, pipeline components, and a model registry.

\begin{center}
\footnotesize
\begin{tabular}{@{}p{3.0cm}p{2.9cm}p{3.1cm}p{3.2cm}@{}}
\toprule
\textbf{Capability} & \textbf{AWS} & \textbf{Azure} & \textbf{GCP} \\
\midrule
ML platform & SageMaker & Azure Machine Learning & Vertex AI \\
AutoML & SageMaker Autopilot & Azure AutoML & Vertex AutoML \\
Pipelines & SageMaker Pipelines & AML pipelines (components) & Vertex Pipelines (Kubeflow) \\
Feature store & SageMaker Feature Store & AML managed feature store & Vertex AI Feature Store \\
Model registry & SageMaker Model Registry & AML Model Registry & Vertex Model Registry \\
\addlinespace
\multicolumn{4}{@{}l@{}}{\textbf{Fabric}: built-in ML with MLflow on lakehouses; \textbf{Snowflake}: Snowpark ML;} \\
\multicolumn{4}{@{}l@{}}{\textbf{MongoDB}: partner MLOps over Atlas; \textbf{Databricks}: managed MLflow + UC model registry.} \\
\bottomrule
\end{tabular}
\end{center}

\noindent MLflow is the lingua franca underneath: AML, Databricks, and Fabric all track experiments
with it, so experiment-tracking skills transfer even where the platform does not. Databricks adds
alias-based lifecycle (champion/challenger) and registry lineage from training data to deployed
model inside Unity Catalog.

\section{Model Serving, Monitoring, and Retraining}
\index{model serving!equivalence}\index{drift monitoring!equivalence}%
Chapter~19's serving splits everywhere into real-time endpoints and batch scoring, and monitoring
splits into infrastructure health, data drift, and model quality.

\begin{center}
\footnotesize
\begin{tabular}{@{}p{3.0cm}p{2.9cm}p{3.1cm}p{3.2cm}@{}}
\toprule
\textbf{Capability} & \textbf{AWS} & \textbf{Azure} & \textbf{GCP} \\
\midrule
Real-time serving & SageMaker endpoints & AML managed online endpoints & Vertex AI endpoints \\
Batch scoring & SageMaker batch transform & AML batch endpoints & Vertex batch prediction \\
Model monitoring & SageMaker Model Monitor & AML model monitoring / drift & Vertex Model Monitoring \\
Registry stages & Registry + approval states & Registry + tags/stages & Registry + aliases \\
Retrain trigger & EventBridge $\rightarrow$ pipeline & Event Grid $\rightarrow$ AML pipeline & Eventarc $\rightarrow$ pipeline \\
\addlinespace
\multicolumn{4}{@{}l@{}}{\textbf{Fabric}: MLflow models + SynapseML serving; \textbf{Snowflake}: Snowpark ML registry;} \\
\multicolumn{4}{@{}l@{}}{\textbf{MongoDB}: in-app scoring over Atlas; \textbf{Databricks}: Model Serving + inference tables + Lakehouse Monitoring.} \\
\bottomrule
\end{tabular}
\end{center}

\noindent The MLOps loop---monitor, detect drift, retrain, gate, promote---is platform-neutral.
The Databricks distinction: inference logs are Delta tables in your own catalog, so monitoring is
a SQL/BI problem over data you already govern, not a separate observability silo.

\section{Unstructured Data and Generative AI}
\index{generative AI!equivalence}\index{vector search!equivalence}%
Chapter~20's document AI, search, and managed foundation models exist on every cloud; the pipeline
pattern---ingest, extract, index, retrieve, generate---is identical across vendors.

\begin{center}
\footnotesize
\begin{tabular}{@{}p{3.0cm}p{2.9cm}p{3.1cm}p{3.2cm}@{}}
\toprule
\textbf{Capability} & \textbf{AWS} & \textbf{Azure} & \textbf{GCP} \\
\midrule
Document extraction & Textract & AI Document Intelligence & Document AI \\
NLP / PII detection & Comprehend & AI Language & Cloud Natural Language \\
Search / vectors & OpenSearch / Kendra & AI Search (vector + hybrid) & Vertex AI Search \\
Foundation models & Bedrock & Azure OpenAI Service & Vertex AI (Gemini) \\
Event-driven processing & Lambda & Azure Functions & Cloud Functions \\
\addlinespace
\multicolumn{4}{@{}l@{}}{\textbf{Fabric}: AI skills over OneLake; \textbf{Snowflake}: Cortex (LLM functions in SQL);} \\
\multicolumn{4}{@{}l@{}}{\textbf{MongoDB}: Atlas Vector Search; \textbf{Databricks}: Foundation Model APIs + Mosaic AI Vector Search.} \\
\bottomrule
\end{tabular}
\end{center}

\noindent The RAG pattern taught in Chapter~20---chunk, embed, index, retrieve, ground the
model---is the same architecture whether the vector store is AI Search, OpenSearch, pgvector,
Atlas Vector Search, or Mosaic AI Vector Search; evaluation with LLM judges (MLflow evaluate,
Bedrock evaluations, Azure AI eval) is likewise a shared discipline.

\section{Zero-Trust Networking, Identity, and Secrets}
\index{zero trust!equivalence}\index{managed identity!equivalence}%
Chapter~21's architecture---private connectivity, no long-lived secrets, centralized keys---is
cloud-universal; the product names differ, the design does not.

\begin{center}
\footnotesize
\begin{tabular}{@{}p{3.0cm}p{2.9cm}p{3.1cm}p{3.2cm}@{}}
\toprule
\textbf{Capability} & \textbf{AWS} & \textbf{Azure} & \textbf{GCP} \\
\midrule
Identity plane & IAM & Microsoft Entra ID + Azure RBAC & Cloud IAM \\
Workload identity & IAM roles for EC2/EKS & Managed identities & Service account impersonation \\
Private service access & PrivateLink / VPC endpoints & Private Endpoints & Private Service Connect \\
Secrets manager & Secrets Manager & Azure Key Vault & Secret Manager \\
Key custody (CMK) & KMS & Key Vault / Managed HSM & Cloud KMS / HSM \\
\addlinespace
\multicolumn{4}{@{}l@{}}{\textbf{Fabric}: Entra ID + OneLake security; \textbf{Snowflake}: PrivateLink + Tri-Secret;} \\
\multicolumn{4}{@{}l@{}}{\textbf{MongoDB}: Atlas private endpoints + CSFLE; \textbf{Databricks}: secret scopes, OAuth M2M, private connectivity.} \\
\bottomrule
\end{tabular}
\end{center}

\noindent The exam and interview favorite---``why is a managed identity better than a service
principal secret?''---has the same answer on every cloud: no secret material exists to leak,
rotate, or commit to git. Databricks' PAT tokens are the cautionary contrast: they are
long-lived secrets, which is why cluster policies and OAuth machine-to-machine service
principals are the production pattern.

\section{Observability, FinOps, and Quality Gates}
\index{observability!equivalence}\index{FinOps!equivalence}%
Chapter~22's operating layer---catalogs, lineage, quality gates, and cost observability---pairs a
native catalog with the same open quality tooling on every cloud.

\begin{center}
\footnotesize
\begin{tabular}{@{}p{3.0cm}p{2.9cm}p{3.1cm}p{3.2cm}@{}}
\toprule
\textbf{Capability} & \textbf{AWS} & \textbf{Azure} & \textbf{GCP} \\
\midrule
Catalog and governance & Glue Data Catalog / Lake Formation & Microsoft Purview & Dataplex \\
Data quality tests & dbt / Great Expectations / Deequ & dbt / Great Expectations in ADF or Databricks & dbt / Great Expectations / Dataplex rules \\
Monitoring & CloudWatch & Azure Monitor + Log Analytics & Cloud Operations \\
Cost management & Cost Explorer + budgets & Cost Management + budgets & Cloud Billing + budgets \\
Sensitivity labels & Macie findings & Purview + MIP labels to Power BI & DLP inspection \\
\addlinespace
\multicolumn{4}{@{}l@{}}{\textbf{Fabric}: Purview hub + Capacity Metrics; \textbf{Snowflake}: Horizon + resource monitors;} \\
\multicolumn{4}{@{}l@{}}{\textbf{MongoDB}: Atlas auditing; \textbf{Databricks}: \texttt{system.billing.usage} + \texttt{system.access.audit} + Lakehouse Monitoring.} \\
\bottomrule
\end{tabular}
\end{center}

\noindent dbt and Great Expectations are deliberately platform-neutral: the same \texttt{dbt test}
or GE checkpoint runs against Synapse, Snowflake, BigQuery, Redshift, or Databricks SQL---quality
tooling is the one layer you can learn once \cite{dbtlabs2025core,greatexpectations2025}.

\section{Data Privacy and Compliance Primitives}
\index{data privacy!equivalence}\index{masking!equivalence}%
Chapter~23's privacy primitives---masking, client-side encryption, erasure, and policy
audit---exist on every platform, and certification exams test them as a family.

\begin{center}
\footnotesize
\begin{tabular}{@{}p{3.0cm}p{2.9cm}p{3.1cm}p{3.2cm}@{}}
\toprule
\textbf{Capability} & \textbf{AWS} & \textbf{Azure} & \textbf{GCP} \\
\midrule
Dynamic masking & Redshift dynamic masking & SQL/Synapse Dynamic Data Masking & BigQuery column-level + masking \\
Client-side encryption & Client-side encryption SDKs & Always Encrypted & Client-side encryption / CSEK \\
At-rest encryption & KMS-backed TDE equivalents & TDE (+ CMK via Key Vault) & Google-managed / CMEK \\
Policy enforcement & AWS Config + SCPs & Azure Policy & Organization Policy \\
Threat detection & GuardDuty / Security Hub & Microsoft Defender for Cloud & Security Command Center \\
PII discovery & Macie & Purview classifications & Sensitive Data Protection \\
\addlinespace
\multicolumn{4}{@{}l@{}}{\textbf{Fabric}: OneLake security + MIP labels; \textbf{Snowflake}: masking policies + tags;} \\
\multicolumn{4}{@{}l@{}}{\textbf{MongoDB}: client-side field-level encryption; \textbf{Databricks}: UC row filters + column masks.} \\
\bottomrule
\end{tabular}
\end{center}

\noindent Crypto-shredding---encrypt per data subject, erase by destroying the subject's key---is
the cross-platform answer to ``delete everywhere, including backups,'' because backups become
unreadable ciphertext when the key is gone.

\section{DataOps: Infrastructure as Code and CI/CD}
\index{infrastructure as code!equivalence}\index{CI/CD!equivalence}%
Chapter~24's disciplines---declarative infrastructure, SQL-based transformation, reviewed
pipelines---are the most portable skills in this book. The same Terraform workflow and the same
dbt project structure run against every cloud.

\begin{center}
\footnotesize
\begin{tabular}{@{}p{3.0cm}p{2.9cm}p{3.1cm}p{3.2cm}@{}}
\toprule
\textbf{Capability} & \textbf{AWS} & \textbf{Azure} & \textbf{GCP} \\
\midrule
Multi-cloud IaC & Terraform \texttt{aws} provider & Terraform \texttt{azurerm} provider & Terraform \texttt{google} provider \\
Native IaC & CloudFormation / CDK & Bicep / ARM templates & Deployment Manager (legacy) \\
Transform-as-code & dbt on Redshift/Athena & dbt on Synapse/Fabric & dbt on BigQuery \\
CI/CD & CodePipeline / GitHub Actions & GitHub Actions / Azure DevOps & Cloud Build / GitHub Actions \\
State backend & S3 + DynamoDB lock & Azure Storage + lease lock & GCS backend \\
\addlinespace
\multicolumn{4}{@{}l@{}}{\textbf{Fabric}: Git integration + deployment pipelines; \textbf{Snowflake}: dbt + Terraform provider;} \\
\multicolumn{4}{@{}l@{}}{\textbf{MongoDB}: Atlas Terraform provider; \textbf{Databricks}: Asset Bundles (\texttt{databricks.yml}) + Terraform provider.} \\
\bottomrule
\end{tabular}
\end{center}

\noindent The lesson that transfers everywhere: \emph{the merge is the deploy, and the pipeline is
the only writer to production}. Databricks Asset Bundles add a validate-before-deploy step
(\texttt{databricks bundle validate}) that mirrors \texttt{terraform plan} and Bicep what-if---the
same reviewed-plan discipline in native tooling.

\begin{tipbox}
Use this appendix as a study aid for the book's lockstep, cross-cloud approach: when you learn a
concept on one platform, immediately locate its equivalent here and predict how the other six
would express the same idea. That habit turns seven separate vocabularies into one mental model.
\end{tipbox}
